% --- Archon physics formalization source begin ---
% archon:physics
% archon:covers IPhO2026Problems/problem_ipho_2026_e1_a1.lean
% archon:source-report reports/ipho_2026/problem_ipho_2026_e1_a1.source.json
% archon:problem-id ipho_2026_e1
% archon:part-id E1-A1

\chapter{Ipho Physics problem ipho\_2026\_e1\_a1}
\label{ch:IPhO2026Problems_problem_ipho_2026_e1_a1}

\paragraph{Problem source.}
\#\# Physical scenario

The experimental apparatus contains a sealed air column (CA) in the inner
cylinder.  Propylene glycol is introduced to h = 4.5 cm so the air volume is
fixed.  Use the cylinder dimensions in Figure 17, ambient air density
rho\_a = 1.12 kg/m\textasciicircum{}3, and the ideal-gas law P*V = n*R*T.  The outer-cylinder
water bath is heated while pressure and temperature are recorded.

\#\# Current subquestion E1-A1

Determine the mass m, amount n, and number N of molecules in the confined air column.

\paragraph{Shared problem context.}
The experimental apparatus contains a sealed air column (CA) in the inner
cylinder.  Propylene glycol is introduced to h = 4.5 cm so the air volume is
fixed.  Use the cylinder dimensions in Figure 17, ambient air density
rho\_a = 1.12 kg/m\textasciicircum{}3, and the ideal-gas law P*V = n*R*T.  The outer-cylinder
water bath is heated while pressure and temperature are recorded.

\paragraph{Current subquestion.}
Determine the mass m, amount n, and number N of molecules in the confined air column.

\paragraph{Figure/image path.}
ipho\_2026\_source/image/E1\_page-9.png

\paragraph{Formalization target.}
create a compiling Lean file with sorry bodies at `IPhO2026Problems/problem\_ipho\_2026\_e1\_a1.lean`.
The Lean declarations must preserve the physical quantities, dimensions or dimensional roles, figure labels, governing-law hypotheses, and final relation expressed by this problem.
Use Mathlib/Physlib names found through LeanExplore where available. If a domain API is missing, introduce faithful local abstractions rather than scalar placeholder aliases.
This is an answer-blind solve-phase task. Derive a candidate result only from the problem-side evidence above; do not assume a target value, select a tolerance around a desired value, or encode a desired result into a candidate domain.
For numerical questions, define the raw end-to-end quantity and apply an explicit source-derived rounding rule. For identification questions, characterize or prove uniqueness before recording the derived candidate.

\begin{theorem}[Ipho Physics formalization target]
\label{thm:physics:ipho_2026_e1_a1:target}
The assigned autoformalize agent should translate this IPhO physics problem into Lean declarations in the covered file, with theorem and lemma proof bodies written as `by sorry`.
\end{theorem}
\begin{proof}
This is an autoformalization task, not a proof task. Produce faithful statements that can later be proved without weakening the source contract.
\end{proof}
% --- Archon physics formalization source end ---
